% Generated by the blueprint tool (blueprint-pdf/1) from a blueprint snapshot; edit the
% Markdown sources, not this file.  Compile with `latexmk -pdf` or by running
% `pdflatex` twice (the second run settles the contents and the page references).
% Images, if any, are read from assets/ next to this file.
\documentclass[11pt]{article}
\usepackage[T1]{fontenc}
\usepackage[utf8]{inputenc}
\usepackage{lmodern}
\usepackage{textcomp}
\usepackage[a4paper,margin=2.5cm,headheight=14pt]{geometry}
\usepackage{amsmath,amssymb,amsthm}
\usepackage{amscd}  % KaTeX's CD diagrams; part of amsmath, so always there
\usepackage{mathtools}
\usepackage{bm}
\usepackage{graphicx}
\usepackage{xcolor}
\usepackage{array,booktabs,tabularx}
\usepackage{fancyvrb}
\usepackage[normalem]{ulem}
\usepackage{microtype}
% Optional, and only with their fonts: a small TeX installation has the
% packages without the fonts, which fails at the first formula.
\IfFileExists{mathrsfs.sty}{\IfFileExists{rsfs10.tfm}{\usepackage{mathrsfs}}{}}{}
\IfFileExists{stmaryrd.sty}{\IfFileExists{stmary10.tfm}{\usepackage{stmaryrd}}{}}{}
\IfFileExists{cancel.sty}{\usepackage{cancel}}{}
\usepackage[hyphens]{url}
\usepackage[pdfusetitle,bookmarksnumbered=false]{hyperref}
\usepackage{bookmark}
\definecolor{bplink}{HTML}{1F4E9C}
\definecolor{bpmuted}{HTML}{5F676E}
\hypersetup{colorlinks=true,linkcolor=bplink,urlcolor=bplink,citecolor=bplink}
\setlength{\parindent}{0pt}
\setlength{\parskip}{0.55em plus 0.2em minus 0.1em}
% Lists as tight as the page's.
\makeatletter
\def\@listi{\leftmargin\leftmargini\parsep 0pt\topsep 0.3em plus 0.1em\itemsep 0.1em plus 0.1em}
\let\@listI\@listi
\def\@listii{\leftmargin\leftmarginii\labelwidth\leftmarginii\advance\labelwidth-\labelsep\parsep 0pt\topsep 0.1em\itemsep 0.1em}
% An ordered list that does not start at 1.
\newcommand{\bpstartat}[1]{\setcounter{\@enumctr}{\numexpr#1-1\relax}}
\makeatother
\setcounter{tocdepth}{3}
\emergencystretch=3em
\allowdisplaybreaks

% A numbered statement, set as a paper sets one: "Theorem 2.1 (Title)."
% #1 is the lead ("Theorem~2.1"), #2 what follows it in roman.
\newenvironment{bpstatement}[2]{\par\addvspace{0.6\baselineskip}\noindent\textbf{#1}#2\textbf{.}\enspace\ignorespaces}
  {\par\addvspace{0.4\baselineskip}}
\newenvironment{bpstatementit}[2]{\begin{bpstatement}{#1}{#2}\itshape\ignorespaces}{\end{bpstatement}}
% A step: a relation with prose of its own, read under the object it starts from.
\newenvironment{bpstep}[1]{\par\addvspace{0.3\baselineskip}\begin{list}{}{\leftmargin=1.6em\rightmargin=0pt\topsep=0pt}\item\relax
  {\small\color{bpmuted}#1}\par\nopagebreak}{\end{list}\par}
% Notes about where an object is written out.
\newenvironment{bpnote}{\par\small\color{bpmuted}\itshape}{\par}
\newcommand{\bpomitted}[1]{\par{\small\color{bpmuted}\itshape #1\par}}
\newcommand{\bpbodyheading}[1]{\par\addvspace{0.4\baselineskip}\noindent\textbf{#1}\par\nopagebreak}
\newcommand{\bprule}{\par\medskip\noindent\hfil\rule{0.3\linewidth}{0.4pt}\hfil\par\medskip}
\newcommand{\bpboundary}[1]{\par\nopagebreak{\small\color{bpmuted}#1\par}\nopagebreak}
% A reference to something this export leaves out, and one that resolves nowhere.
\newcommand{\bpout}[1]{#1\textsuperscript{\hyperlink{bp-outside}{\ensuremath{\ast}}}}
\newcommand{\bpbroken}[1]{\textcolor{red!70!black}{[#1?]}}
\newcommand{\bpslug}[1]{\texttt{#1}}
% A bibliography entry: label, anchor, entry; the label hangs in the margin.
\newcommand{\bpbibitem}[3]{\par\setbox0\hbox{[#1]\enspace}\ifdim\wd0<5em\setbox0\hbox to 5em{[#1]\hfil}\fi
  \hangindent=\wd0 \hangafter=1 \noindent\box0\phantomsection\label{#2}#3\par\smallskip}
% A character the export knows no way to set.
\newcommand{\bpglyph}[1]{\fbox{\scriptsize #1}}
% Images, at their own size unless they are wider than the line.
\newsavebox{\bpimagebox}
\newcommand{\bpimage}[1]{\sbox{\bpimagebox}{\includegraphics{#1}}%
  \ifdim\wd\bpimagebox>\linewidth\resizebox{\linewidth}{!}{\usebox{\bpimagebox}}\else\usebox{\bpimagebox}\fi}
\newcommand{\bpmissingimage}[1]{\fbox{\parbox{0.8\linewidth}{\centering\small\color{bpmuted}Image not included: #1}}}
\DefineVerbatimEnvironment{bpverbatim}{Verbatim}{fontsize=\small,xleftmargin=1em}
\renewcommand{\qedsymbol}{\ensuremath{\blacksquare}}
\hypersetup{pdftitle={Weil Conjectures — D Cohomology of torsion sheaves on curves},pdfsubject={Appendix D Cohomology of torsion sheaves on curves; all levels},pdfcreator={blueprint}}
% KaTeX commands the website knows
\providecommand{\R}{\mathbb{R}}\providecommand{\N}{\mathbb{N}}\providecommand{\Z}{\mathbb{Z}}
\providecommand{\Q}{\mathbb{Q}}\providecommand{\C}{\mathbb{C}}
\providecommand{\Reals}{\mathbb{R}}\providecommand{\reals}{\mathbb{R}}\providecommand{\Complex}{\mathbb{C}}
\providecommand{\cnums}{\mathbb{C}}\providecommand{\natnums}{\mathbb{N}}\providecommand{\Bbb}[1]{\mathbb{#1}}
\providecommand{\bold}[1]{\mathbf{#1}}\providecommand{\frak}[1]{\mathfrak{#1}}\providecommand{\mathbbm}[1]{\mathbb{#1}}
\providecommand{\lang}{\langle}\providecommand{\rang}{\rangle}\providecommand{\bra}[1]{\langle #1\rvert}
\providecommand{\ket}[1]{\lvert #1\rangle}\providecommand{\braket}[1]{\langle #1\rangle}
\providecommand{\set}[1]{\{#1\}}\providecommand{\Set}[1]{\left\{#1\right\}}
\providecommand{\infin}{\infty}\providecommand{\isin}{\in}\providecommand{\plusmn}{\pm}\providecommand{\sub}{\subset}
\providecommand{\sube}{\subseteq}\providecommand{\supe}{\supseteq}\providecommand{\exist}{\exists}
\providecommand{\bull}{\bullet}\providecommand{\larr}{\leftarrow}\providecommand{\rarr}{\rightarrow}
\providecommand{\lrarr}{\leftrightarrow}\providecommand{\Larr}{\Leftarrow}\providecommand{\Rarr}{\Rightarrow}
\providecommand{\Lrarr}{\Leftrightarrow}\providecommand{\uarr}{\uparrow}\providecommand{\darr}{\downarrow}
\providecommand{\harr}{\leftrightarrow}\providecommand{\lArr}{\Leftarrow}\providecommand{\rArr}{\Rightarrow}
\providecommand{\hArr}{\Leftrightarrow}\providecommand{\alef}{\aleph}\providecommand{\alefsym}{\aleph}
\providecommand{\real}{\Re}\providecommand{\image}{\Im}\providecommand{\weierp}{\wp}
\providecommand{\Alpha}{\mathrm{A}}\providecommand{\Beta}{\mathrm{B}}\providecommand{\Epsilon}{\mathrm{E}}
\providecommand{\Zeta}{\mathrm{Z}}\providecommand{\Eta}{\mathrm{H}}\providecommand{\Iota}{\mathrm{I}}
\providecommand{\Kappa}{\mathrm{K}}\providecommand{\Mu}{\mathrm{M}}\providecommand{\Nu}{\mathrm{N}}
\providecommand{\Omicron}{\mathrm{O}}\providecommand{\Rho}{\mathrm{P}}\providecommand{\Tau}{\mathrm{T}}
\providecommand{\Chi}{\mathrm{X}}\providecommand{\omicron}{o}\providecommand{\gt}{>}\providecommand{\lt}{<}
\providecommand{\And}{\&}\providecommand{\hdashline}{\hline}\providecommand{\operatornamewithlimits}[1]{\operatorname*{#1}}
\providecommand{\llbracket}{[\![}\providecommand{\rrbracket}{]\!]}\providecommand{\cancel}[1]{#1}
\providecommand{\mathscr}{\mathcal}
% the characters beyond ASCII this document uses
\DeclareUnicodeCharacter{2160}{\textup{I}}
% the project's macros ([katex.macros] in blueprint.toml)
\providecommand{\A}{}\renewcommand{\A}{\mathbf A}
\providecommand{\AEn}{}\renewcommand{\AEn}{A_{E,n}}
\providecommand{\Aaf}{}\renewcommand{\Aaf}{\mathbf A}
\providecommand{\Ab}{}\renewcommand{\Ab}{\mathsf{Ab}}
\providecommand{\AffSch}{}\renewcommand{\AffSch}{\operatorname{AffSch}}
\providecommand{\Aone}{}\renewcommand{\Aone}{\mathbf A^1}
\providecommand{\Bl}{}\renewcommand{\Bl}{\operatorname{Bl}}
\providecommand{\Br}{}\renewcommand{\Br}{\mathrm{Br}}
\providecommand{\C}{}\renewcommand{\C}{\mathbb C}
\providecommand{\Cech}{}\renewcommand{\Cech}{\check{H}}
\providecommand{\Coext}{}\renewcommand{\Coext}{\operatorname{Coext}}
\providecommand{\Cone}{}\renewcommand{\Cone}{\operatorname{Cone}}
\providecommand{\Cons}{}\renewcommand{\Cons}{\operatorname{Cons}}
\providecommand{\Cref}{}\renewcommand{\Cref}[1]{\text{#1}}
\providecommand{\D}{}\renewcommand{\D}{D}
\providecommand{\DLAFT}{}\renewcommand{\DLAFT}{\operatorname{FT}}
\providecommand{\DLAFr}{}\renewcommand{\DLAFr}{\operatorname{Fr}}
\providecommand{\DLAQ}{}\renewcommand{\DLAQ}{{\overline{\mathbf Q}_\ell}}
\providecommand{\DLRQ}{}\renewcommand{\DLRQ}{{\overline{\mathbf Q}_\ell}}
\providecommand{\Dbc}{}\renewcommand{\Dbc}{D^b_c}
\providecommand{\Dc}{}\renewcommand{\Dc}{D_{\mathrm{cons}}}
\providecommand{\Dctf}{}\renewcommand{\Dctf}{D^b_{ctf}}
\providecommand{\Dic}{}\renewcommand{\Dic}{D_{\IndCons}}
\providecommand{\Dplus}{}\renewcommand{\Dplus}{D^{+}}
\providecommand{\End}{}\renewcommand{\End}{\operatorname{End}}
\providecommand{\Et}{}\renewcommand{\Et}{\mathrm{\acute{e}t}}
\providecommand{\Ex}{}\renewcommand{\Ex}{\mathrm{Ex}}
\providecommand{\Ext}{}\renewcommand{\Ext}{\operatorname{Ext}}
\providecommand{\F}{}\renewcommand{\F}{\mathbb F}
\providecommand{\FT}{}\renewcommand{\FT}{\operatorname{FT}}
\providecommand{\Fbar}{}\renewcommand{\Fbar}{{\overline{\mathbf F}_q}}
\providecommand{\Fell}{}\renewcommand{\Fell}{\mathbf F_\ell}
\providecommand{\Fib}{}\renewcommand{\Fib}{\operatorname{Fib}}
\providecommand{\Fix}{}\renewcommand{\Fix}{\operatorname{Fix}}
\providecommand{\Fp}{}\renewcommand{\Fp}{{\mathbf F_{p}}}
\providecommand{\Fq}{}\renewcommand{\Fq}{{\mathbf F_q}}
\providecommand{\Fqn}{}\renewcommand{\Fqn}{{\mathbf F_{q^{n}}}}
\providecommand{\Frob}{}\renewcommand{\Frob}{\operatorname{Frob}}
\providecommand{\GL}{}\renewcommand{\GL}{\operatorname{GL}}
\providecommand{\Gal}{}\renewcommand{\Gal}{\operatorname{Gal}}
\providecommand{\Hens}{}\renewcommand{\Hens}{\operatorname{Hens}}
\providecommand{\Hilb}{}\renewcommand{\Hilb}{\operatorname{Hilb}}
\providecommand{\Hom}{}\renewcommand{\Hom}{\operatorname{Hom}}
\providecommand{\IndCons}{}\renewcommand{\IndCons}{\operatorname{IndCons}}
\providecommand{\Kar}{}\renewcommand{\Kar}{\operatorname{Kar}}
\providecommand{\LT}{}\renewcommand{\LT}{\operatorname{LT}}
\providecommand{\Lie}{}\renewcommand{\Lie}{\operatorname{Lie}}
\providecommand{\Loc}{}\renewcommand{\Loc}{\operatorname{Loc}}
\providecommand{\Mod}{}\renewcommand{\Mod}{\mathsf{Mod}}
\providecommand{\N}{}\renewcommand{\N}{\mathbb{N}}
\providecommand{\Needspace}{}\renewcommand{\Needspace}{\par}
\providecommand{\O}{}\renewcommand{\O}{{\mathcal{O}}}
\providecommand{\Ob}{}\renewcommand{\Ob}{\mathrm{Ob}}
\providecommand{\Oe}{}\renewcommand{\Oe}{\mathcal O_E}
\providecommand{\PP}{}\renewcommand{\PP}{\mathbf{P}}
\providecommand{\Perf}{}\renewcommand{\Perf}{\operatorname{Perf}}
\providecommand{\Pic}{}\renewcommand{\Pic}{\mathrm{Pic}}
\providecommand{\Pm}{}\renewcommand{\Pm}{\mathbf P^{m}}
\providecommand{\Pone}{}\renewcommand{\Pone}{\mathbf P^{1}}
\providecommand{\Pqn}{}\renewcommand{\Pqn}{\mathbf P^{n}}
\providecommand{\ProFin}{}\renewcommand{\ProFin}{\operatorname{ProFin}}
\providecommand{\Proet}{}\renewcommand{\Proet}{\mathrm{pro\acute{e}t}}
\providecommand{\Q}{}\renewcommand{\Q}{\mathbb Q}
\providecommand{\Qb}{}\renewcommand{\Qb}{\overline{\mathbf Q}_{\ell}}
\providecommand{\Qbar}{}\renewcommand{\Qbar}{\overline{\mathbf Q}}
\providecommand{\QbarX}{}\renewcommand{\QbarX}{\overline{\mathbf Q}_{\ell,X}}
\providecommand{\Qbarl}{}\renewcommand{\Qbarl}{{\overline{\mathbf Q}_{\ell}}}
\providecommand{\Qell}{}\renewcommand{\Qell}{{\mathbf{Q}_{\ell}}}
\providecommand{\Ql}{}\renewcommand{\Ql}{\mathbf Q_{\ell}}
\providecommand{\Qlb}{}\renewcommand{\Qlb}{{\overline{\mathbf Q}_\ell}}
\providecommand{\Qlbar}{}\renewcommand{\Qlbar}{\overline{\mathbf Q}_\ell}
\providecommand{\R}{}\renewcommand{\R}{\mathbb{R}}
\providecommand{\RG}{}\renewcommand{\RG}{\operatorname{\mathrm{R}\Gamma}}
\providecommand{\RGamma}{}\renewcommand{\RGamma}{R\Gamma}
\providecommand{\RGammac}{}\renewcommand{\RGammac}{R\Gamma_c}
\providecommand{\RHom}{}\renewcommand{\RHom}{R\mathcal{H}om}
\providecommand{\RHomX}{}\renewcommand{\RHomX}{R\mathcal H om_X}
\providecommand{\Rat}{}\renewcommand{\Rat}{\operatorname{IsRat}}
\providecommand{\Res}{}\renewcommand{\Res}{\operatorname{Res}}
\providecommand{\Rlim}{}\renewcommand{\Rlim}{R\!\varprojlim}
\providecommand{\SIm}{}\renewcommand{\SIm}{\operatorname{SIm}}
\providecommand{\Sch}{}\renewcommand{\Sch}{\mathsf{Sch}}
\providecommand{\Sh}{}\renewcommand{\Sh}{\operatorname{Sh}}
\providecommand{\Span}{}\renewcommand{\Span}{\operatorname{span}}
\providecommand{\Spec}{}\renewcommand{\Spec}{\operatorname{Spec}}
\providecommand{\Tr}{}\renewcommand{\Tr}{\operatorname{Tr}}
\providecommand{\Trsm}{}\renewcommand{\Trsm}{\operatorname{Tr}^{\mathrm{sm}}}
\providecommand{\Wt}{}\renewcommand{\Wt}{\operatorname{Wt}}
\providecommand{\Z}{}\renewcommand{\Z}{\mathbb{Z}}
\providecommand{\Zar}{}\renewcommand{\Zar}{\mathrm{zar}}
\providecommand{\Zell}{}\renewcommand{\Zell}{{\mathbf{Z}_{\ell}}}
\providecommand{\Zl}{}\renewcommand{\Zl}{{\mathbf Z_\ell}}
\providecommand{\affet}{}\renewcommand{\affet}[1]{{#1}^{\mathrm{aff}}_{\mathrm{ét}}}
\providecommand{\affproet}{}\renewcommand{\affproet}[1]{{#1}^{\mathrm{aff}}_{\mathrm{proét}}}
\providecommand{\boxt}{}\renewcommand{\boxt}{\boxtimes_A^{\mathrm L}}
\providecommand{\cF}{}\renewcommand{\cF}{\mathcal{F}}
\providecommand{\cHom}{}\renewcommand{\cHom}{R\mathcal{H}\!om}
\providecommand{\cI}{}\renewcommand{\cI}{\mathcal{I}}
\providecommand{\cL}{}\renewcommand{\cL}{\mathcal L}
\providecommand{\cM}{}\renewcommand{\cM}{\mathcal M}
\providecommand{\cN}{}\renewcommand{\cN}{\mathcal N}
\providecommand{\cO}{}\renewcommand{\cO}{\mathcal O}
\providecommand{\cT}{}\renewcommand{\cT}{\mathcal T}
\providecommand{\checked}{}\renewcommand{\checked}{\color{red}\aftergroup\restorecolor}
\providecommand{\coeff}{}\renewcommand{\coeff}{\operatorname{coeff}}
\providecommand{\coker}{}\renewcommand{\coker}{\operatorname{coker}}
\providecommand{\colim}{}\renewcommand{\colim}{\operatorname{colim}}
\providecommand{\constructedby}{}\renewcommand{\constructedby}[1]{\par\smallskip\noindent\textcolor{purple!60!black}{\textsf{Constructed by:} }
  \footnotesize\sloppy\nolinkurl{#1}\par}
\providecommand{\deepresult}{}\renewcommand{\deepresult}[1]{\par\smallskip\noindent\textcolor{orange!45!black}{
  \textsf{Blueprint role: named deep theorem node; see the cluster proof architecture.}}
  \par\noindent\textcolor{orange!45!black}{\textsf{Mathematical source:} }#1\par}
\providecommand{\dt}{}\renewcommand{\dt}{\mathbin{\otimes^{\mathbf L}}}
\providecommand{\dw}{}\renewcommand{\dw}{\operatorname{dw}}
\providecommand{\et}{}\renewcommand{\et}[1]{{#1}_{\mathrm{ét}}}
\providecommand{\ev}{}\renewcommand{\ev}{\operatorname{ev}}
\providecommand{\fib}{}\renewcommand{\fib}{\operatorname{fib}}
\providecommand{\formalizationnote}{}\renewcommand{\formalizationnote}[1]{\begin{quote}\small\raggedright\textbf{Formalization note.} #1\end{quote}}
\providecommand{\fpqc}{}\renewcommand{\fpqc}{X_{\mathrm{fpqc}}}
\providecommand{\gammacont}{}\renewcommand{\gammacont}{\Gamma_{\mathrm{cont}}}
\providecommand{\gm}{}\renewcommand{\gm}{\mathbf G_m}
\providecommand{\hocolim}{}\renewcommand{\hocolim}{\operatorname{hocolim}}
\providecommand{\holim}{}\renewcommand{\holim}{\operatorname{holim}}
\providecommand{\hrB}{}\renewcommand{\hrB}{\@ifnextchar]{\hrEnd}{\hrStep}}
\providecommand{\hrEnd}{}\renewcommand{\hrEnd}{\@firstoftwo{}}
\providecommand{\hrStep}{}\renewcommand{\hrStep}{\@firstoftwo{\hrB}}
\providecommand{\id}{}\renewcommand{\id}{\operatorname{id}}
\providecommand{\kunFl}{}\renewcommand{\kunFl}{\mathbf F_\ell}
\providecommand{\kunZn}{}\renewcommand{\kunZn}{\mathbf Z/n\mathbf Z}
\providecommand{\kuncd}{}\renewcommand{\kuncd}{\operatorname{cd}}
\providecommand{\kundim}{}\renewcommand{\kundim}{\dim}
\providecommand{\kunotimes}{}\renewcommand{\kunotimes}[1]{\otimes^{\mathbf L}_{#1}}
\providecommand{\kunsupp}{}\renewcommand{\kunsupp}{\operatorname{Supp}}
\providecommand{\lotimes}{}\renewcommand{\lotimes}{\otimes^{\mathrm L}}
\providecommand{\m}{}\renewcommand{\m}{{\mathfrak{m}}}
\providecommand{\n}{}\renewcommand{\n}{{\mathfrak{n}}}
\providecommand{\nagataGap}{}\renewcommand{\nagataGap}[1]{\par\noindent\textbf{Formalization gap.} #1\par}
\providecommand{\openobligation}{}\renewcommand{\openobligation}[1]{\par\smallskip\noindent\textcolor{red!65!black}{\textsf{Open obligation:} }
  \footnotesize\sloppy\nolinkurl{#1}\par}
\providecommand{\p}{}\renewcommand{\p}{{\mathfrak{p}}}
\providecommand{\pr}{}\renewcommand{\pr}{\operatorname{pr}}
\providecommand{\preshv}{}\renewcommand{\preshv}[1]{\mathrm{PreShv}(#1)}
\providecommand{\pro}{}\renewcommand{\pro}[1]{\mathrm{Pro}(#1)}
\providecommand{\proet}{}\renewcommand{\proet}{X_{\mathrm{proét}}}
\providecommand{\q}{}\renewcommand{\q}{{\mathfrak{q}}}
\providecommand{\rad}{}\renewcommand{\rad}{\operatorname{rad}}
\providecommand{\rank}{}\renewcommand{\rank}{\operatorname{rank}}
\providecommand{\red}{}\renewcommand{\red}{\color{red}}
\providecommand{\restorecolor}{}\renewcommand{\restorecolor}{\color{black}}
\providecommand{\retiredobligation}{}\renewcommand{\retiredobligation}[1]{\par\smallskip\noindent\textcolor{red!75!black}{\textsf{Retired invalid route:} }
  \footnotesize\sloppy\nolinkurl{#1}\par}
\providecommand{\rk}{}\renewcommand{\rk}{\operatorname{rk}}
\providecommand{\sep}{}\renewcommand{\sep}{\textnormal{sep}}
\providecommand{\shv}{}\renewcommand{\shv}[1]{\mathrm{Shv}(#1)}
\providecommand{\spec}{}\renewcommand{\spec}[1]{\mathrm{Spec}(#1)}
\providecommand{\stacks}{}\renewcommand{\stacks}[1]{\href{https://stacks.math.columbia.edu/tag/#1}{\textnormal{Stacks, Tag~#1}}}
\providecommand{\stacksproject}{}\renewcommand{\stacksproject}[1]{{\cite[\href{https://stacks.math.columbia.edu/tag/#1}{Tag #1}]{stacks-project}}}
\providecommand{\statusnote}{}\renewcommand{\statusnote}[1]{\begin{quote}\small\raggedright\textbf{Status.} #1\end{quote}}
\providecommand{\tildeideal}{}\renewcommand{\tildeideal}[2]{\widetilde{(#1, #2)}}
\providecommand{\tildering}{}\renewcommand{\tildering}[3]{{#1}_{\widetilde{(#2, #3)}}}
\providecommand{\tr}{}\renewcommand{\tr}{\operatorname{tr}}
\providecommand{\wt}{}\renewcommand{\wt}{\operatorname{wt}}

\begin{document}
\pagestyle{plain}
\begin{center}
{\LARGE\bfseries Weil Conjectures\par}
\vspace{0.7em}
{\large Appendix D: Cohomology of torsion sheaves on curves\par}
\vspace{0.5em}
{\small\color{bpmuted}Snapshot 2026-09-26, revision 8e4cb88e2a\par}
\end{center}
\vspace{0.5em}
{\small\tableofcontents}
\bigskip\hrule\bigskip

% object: chap-Cohomology-of-torsion-sheaves-on-curves (blueprint/chap-Cohomology-of-torsion-sheaves-on-curves/_section.md)
\section*{\texorpdfstring{D\quad Cohomology of torsion sheaves on curves}{D Cohomology of torsion sheaves on curves}}
\phantomsection\label{bp1}%
\addcontentsline{toc}{section}{\texorpdfstring{D\quad Cohomology of torsion sheaves on curves}{D Cohomology of torsion sheaves on curves}}
% text: prose-D-01
\phantomsection\label{bp2}%
The last step of proper base change for \(\mathbf P^1\) compares the cohomology of a curve over two separably closed fields, one inside the other. Smooth base change needs the same comparison for affine curves, with coefficients prime to the characteristic. Both rest on Theorem \bpout{4.20}, whose assertions we recall. Let \(k\) be algebraically closed, let \(C\) be a separated scheme of finite type over \(k\) with \(\dim C\leq1\), and let \(F\) be a classical torsion abelian sheaf on \(C_{\mathrm{pro}\acute{e}t}\). Then:

\begin{itemize}
\item \((C_{1})\) \(H^q(C,F)=0\) for \(q>2\);
\item \((C_{2})\) if \(C\) is affine, \(H^q(C,F)=0\) for \(q>1\);
\item \((C_{3})\) if \(\operatorname{char}k=p>0\) and \(F\) is \(p\)-primary, \(H^q(C,F)=0\) for \(q>1\);
\item \((C_{4})\) if \(F\) is constructible with finite stalks and prime to \(\operatorname{char}k\), all \(H^q(C,F)\) are finite;
\item \((C_{5})\) if \(C\) is proper and \(F\) is constructible with finite stalks, all \(H^q(C,F)\) are finite;
\item \((C_{6})\) if \(k'/k\) is an extension of algebraically closed fields and \(F\) is prime to \(\operatorname{char}k\), the pullback maps \(H^q(C,F)\to H^q(C_{k'},F|)\) are isomorphisms;
\item \((C_{7})\) if \(C\) is proper, the maps in \((C_6)\) are isomorphisms for every torsion \(F\);
\item \((C_{8})\) for any open \(U\subseteq C\), restriction \(H^2(C,F)\to H^2(U,F|_U)\) is surjective.
\end{itemize}

Over \(\mathbf C\) the vanishing and finiteness assertions are statements about a Riemann surface, which has the homotopy type of a CW complex of dimension two, or of one if it is not compact.

In algebraic geometry we compute the groups. Let \(C\) be a smooth proper connected curve of genus \(g\) and let \(n\) be invertible in \(k\). Three ingredients compute the cohomology of \(\mu_n\): the Kummer sequence \(0\to\mu_n\to\mathbf G_m\xrightarrow{n}\mathbf G_m\to0\), the vanishing of \(H^q(C,\mathbf G_m)\) for \(q\geq2\), which is a consequence of Tsen's theorem, and the Jacobian of Appendix \bpout{C}. They give

\[
H^0(C,\mu_n)=\mu_n(k),\qquad H^1(C,\mu_n)=J(C)[n](k)\simeq(\mathbf Z/n)^{2g},
\qquad H^2(C,\mu_n)=\mathbf Z/n
\]

(Proposition \hyperref[bp10]{D.3}). These groups do not change under extension of algebraically closed fields, because torsion points of an abelian variety do not. For \(\mathbf P^1\), which has \(g=0\), the computation fixes the normalization of the class of \(\mathcal O(1)\) that Remark \hyperref[bp12]{D.4} provides for the proof of proper base change. Removing a nonempty finite set of \(r\) points kills \(H^2\) and raises the rank of \(H^1\) to \(2g+r-1\) (Proposition \hyperref[bp15]{D.5}).

Proper base change, however, is needed for all torsion coefficients, including those divisible by the characteristic \(p\), and there the answer changes. On the affine line, \(H^1(\mathbf A^1_k,\mathbf F_p)=k[t]/\{h^p-h\}\) is infinite, and it grows when \(k\) is enlarged (Example \hyperref[bp26]{D.8}). Indeed, the Artin–Schreier torsors \(y^p-y=at\), for \(a\in k\), have pairwise distinct classes, and over a larger field \(k'\) there are new ones with \(a\in k'\setminus k\). So for \(p\)-torsion coefficients, finiteness and invariance under field extension hold only for proper curves. This is why \((C_5)\) and \((C_7)\) of Theorem \bpout{4.20} assume properness, and why Lemma \bpout{4.21} is stated for proper curves.

The Artin–Schreier sequence \(0\to\mathbf F_p\to\mathbf G_a\xrightarrow{\varphi-1}\mathbf G_a\to0\) reduces the cohomology of \(\mathbf F_p\) to the coherent cohomology of \(\mathcal O_C\). On the affine line it presents \(H^1\) as the cokernel of \(h\mapsto h^p-h\) on \(k[t]\), as above. On every curve it gives \(H^q(C,\mathbf F_p)=0\) for \(q\geq2\) (Theorem \bpout{4.20}, \((C_3)\)).

A general torsion sheaf is reduced to a constant one by two observations. First, a sheaf supported on finitely many points has no cohomology in positive degrees. Hence passing to the normalization of a curve, or extending a sheaf by zero from a dense open subset, changes cohomology only in degrees zero and one, and only by finite groups (Lemmas \hyperref[bp36]{D.11} and \hyperref[bp41]{Lemma D.12}). Second, a finite \(\mathbf F_\ell\)-local system becomes a successive extension of constant sheaves on a finite étale cover of degree \(d\) prime to \(\ell\). This cover is given by an \(\ell\)-Sylow subgroup of the monodromy. Dividing the trace of the cover by \(d\) shows that the local system is a direct summand of the pushforward of its pullback (Lemma \hyperref[bp46]{D.13}). We follow the proof of [Stacks, \href{https://stacks.math.columbia.edu/tag/03SC}{03SC}].

Throughout this appendix, cohomology means cohomology on \(C_{\mathrm{pro}\acute{e}t}\), and the sheaves are classical in the sense of Definition \bpout{2.30}: by Theorem \bpout{2.32} such a sheaf is the inverse image \(\nu_C^*G\) of a sheaf \(G\) on \(C_{\acute{e}t}\) along the comparison morphism. By Theorem \bpout{2.35}, its cohomology, together with the restriction and pullback maps, agrees with étale cohomology, so we carry out all calculations on \(C_{\acute{e}t}\) and transport them back by \(\nu_C^*\).

Let \(p\ge0\) be the characteristic of \(k\). In this appendix \(\ell\) is an arbitrary prime, possibly equal to \(p\). In characteristic zero the hypotheses in \((C_4)\) and \((C_6)\) of Theorem \bpout{4.20} that \(F\) be prime to the characteristic hold automatically. We begin with direct images along finite morphisms, which all the reductions below use.

% object: lemma-finite-direct-images (blueprint/chap-Cohomology-of-torsion-sheaves-on-curves/lemma-finite-direct-images/_section.md)
\phantomsection\label{bp3}%
\begin{bpstatementit}{Lemma~D.1}{ (Finite direct images)}
Let \(u:V\to W\) be finite. For an abelian sheaf \(F\) on \(V_{\acute{e}t}\) and a geometric point \(\bar w:\operatorname{Spec}(\Omega)\to W\), with \(\Omega\) algebraically closed, one has

\[
(u_*F)_{\bar w}
\simeq \prod_{\substack{\bar v:\operatorname{Spec}(\Omega)\to V\\
u\bar v=\bar w}}F_{\bar v}.
\]

The product is finite. The functor \(u_*\) from abelian sheaves on \(V_{\acute{e}t}\) to abelian sheaves on \(W_{\acute{e}t}\) is exact. The same holds for sheaves of modules over a constant coefficient ring.
\end{bpstatementit}

% proof: proof/lemma-finite-direct-images
\begin{proof}[{Proof}]\phantomsection\label{bp4}%
Over a strict localization the source splits into strictly henselian local pieces, and on such a piece taking sections is the same as taking a stalk.

After base change to the strict henselization of \(W\) at \(\bar w\), the finite scheme \(V\) becomes a finite disjoint union of strictly henselian local schemes. On each such local scheme, the sections of an étale sheaf are its stalk at the closed geometric point. Indeed, every étale neighbourhood of that point admits a section, so a germ extends to a global section, and two sections with the same germ agree after pulling back along such a section. In particular, taking sections is exact on abelian sheaves. Combining this with the description of stalks of direct images via strict localizations [Stacks, \href{https://stacks.math.columbia.edu/tag/03QP}{03QP}], we obtain the displayed formula. Since geometric stalks are conservative and finite products of abelian groups are exact, \(u_*\) is exact. The same argument applies to sheaves of modules.
\end{proof}

% object: sec-Kummer-theory (blueprint/chap-Cohomology-of-torsion-sheaves-on-curves/sec-Kummer-theory/_section.md), content hidden
\subsection*{\texorpdfstring{D.1\quad Kummer theory}{D.1 Kummer theory}}
\phantomsection\label{bp5}%
\addcontentsline{toc}{subsection}{\texorpdfstring{D.1\quad Kummer theory}{D.1 Kummer theory}}
% text: prose-D-02
\phantomsection\label{bp6}%
Let \(C\) be a smooth curve over an algebraically closed field \(k\), and let \(\nu_C^*\mathbf G_m\) be the pullback of the multiplicative group \(\mathbf G_m\) of \(C_{\acute{e}t}\) to \(C_{\mathrm{pro}\acute{e}t}\). For \(n\) invertible in \(k\), the Kummer sequence \(0\to\mu_n\to\mathbf G_m\xrightarrow{n}\mathbf G_m\to0\) reduces the cohomology of \(\mu_n\) to that of \(\nu_C^*\mathbf G_m\), so we compute the latter first. In degree one, \(H^1(C,\mathbf G_m)\) is the Picard group, as in the Zariski topology, while the groups \(H^q(C,\mathbf G_m)\) vanish for \(q\ge2\) (Lemma \hyperref[bp7]{D.2}). In particular the Brauer group of \(C\) is zero. This matches, over \(\mathbf C\), the topology of a Riemann surface, whose integral cohomology vanishes above degree two, and whose \(H^2\) is accounted for by the degree of line bundles. In algebraic geometry the vanishing comes from Tsen's theorem, which says that forms of low degree over the function field \(k(C)\) have nontrivial zeros. Over a finite field Tsen's theorem fails: by class field theory the function field of a curve has nonzero Brauer group. This is one reason why we work over an algebraically closed field. With the lemma and the Jacobian, the Kummer sequence then computes the cohomology of \(\mu_n\) on a proper curve (Proposition \hyperref[bp10]{D.3}).

% object: lemma-cohomology-multiplicative-group (blueprint/chap-Cohomology-of-torsion-sheaves-on-curves/sec-Kummer-theory/lemma-cohomology-multiplicative-group/_section.md)
\phantomsection\label{bp7}%
\begin{bpstatementit}{Lemma~D.2}{ (Cohomology of the multiplicative group)}
For a smooth separated curve \(C\) of finite type over an algebraically closed field, one has

\[
H^1(C,\nu_C^*\mathbf G_m)=\operatorname{Pic}(C),
\qquad
H^q(C,\nu_C^*\mathbf G_m)=0\quad(q\geq2).
\]

The identifications are natural under restriction and extension of the algebraically closed ground field. In particular \(\operatorname{Br}(C)=H^2(C_{\acute{e}t},\mathbf G_m)=0\).
\end{bpstatementit}

% proof: proof/lemma-cohomology-multiplicative-group
\begin{proof}[{Proof}]\phantomsection\label{bp8}%
We compute on \(C_{\acute{e}t}\). In degree one, descent for invertible modules identifies \(\mathbf G_m\)-torsors with line bundles. In degrees at least two, by contrast, the vanishing comes from Tsen's theorem, which we prove first (see also [Stacks, \href{https://stacks.math.columbia.edu/tag/03RM}{03RM}]). There are three steps: Tsen's theorem, proved by counting dimensions with Riemann–Roch; the vanishing of the Brauer group and of the Galois cohomology of \(k(C)\); and the passage from the generic point to \(C\) by the divisor sequence.

Tsen's theorem says that \(k(C)\) is \(C_1\), that is, a homogeneous form of degree \(d\) in \(m>d\) variables with coefficients in \(k(C)\) has a nonzero zero.  To see this, let \(\overline C\) be the smooth proper model of \(k(C)\), of genus \(g\), choose an effective divisor \(D\) on \(\overline C\) bounding the poles of all coefficients, and let the unknowns range over \(H^0(\overline C,\mathcal O(NH))\) for an effective divisor \(H\) of positive degree.  The value of the form then lies in \(H^0(\overline C,\mathcal O(dNH+D))\).  Once \(N\deg H>2g-2\), we know both dimensions by Theorem \bpout{C.3}.  In terms of bases, the \(m\) unknowns amount to \(m(N\deg H+1-g)\) unknowns in \(k\). Each coordinate of the value, in a basis of the target space, is a homogeneous polynomial of degree \(d\) in these unknowns, so the vanishing of the value amounts to \(dN\deg H+\deg D+1-g\) homogeneous polynomial equations in them.  The number of unknowns minus the number of equations is \((m-d)N\deg H+(m-1)(1-g)-\deg D\), which is positive for large \(N\) because \(m>d\).  Since \(k\) is algebraically closed, homogeneous equations in more unknowns than equations have a nonzero solution [Stacks, \href{https://stacks.math.columbia.edu/tag/03RD}{03RD}].

The reduced norm form of a division algebra of degree \(d>1\) over \(k(C)\) has degree \(d\) in \(d^2>d\) variables and does not vanish on a nonzero element.  Therefore Tsen's theorem implies that \(\operatorname{Br}(k(C))=0\) [Stacks, \href{https://stacks.math.columbia.edu/tag/03RF}{03RF}].  The same holds for every finite extension of \(k(C)\), since it is again the function field of a curve over \(k\). This is enough for Galois cohomology, because a field all of whose finite extensions have trivial Brauer group has \(H^q(K,(K^{\mathrm{sep}})^\times)=0\) for all \(q\geq1\) [Stacks, \href{https://stacks.math.columbia.edu/tag/03R8}{03R8}]; hence this holds for \(K=k(C)\). To pass from \(k(C)\) to \(C\) we use the divisor sequence, which relates \(\mathbf G_m\) on \(C_{\acute{e}t}\) to the pushforward of the multiplicative group at the generic point and to skyscraper sheaves \(\mathbf Z\) at the closed points. Of these, the former has no higher cohomology by this Galois vanishing, and the latter have none because the residue fields are algebraically closed. Therefore, from the long exact cohomology sequence, we conclude that \(H^q(C_{\acute{e}t},\mathbf G_m)=0\) for \(q\geq2\).

By Theorem \bpout{2.35} these identifications carry over to \(C_{\mathrm{pro}\acute{e}t}\), and they are natural because the comparison maps, the identification of torsors and the restriction maps are functorial.
\end{proof}

% text: prose-D-03
\phantomsection\label{bp9}%
On a proper curve, the lemma and the Kummer sequence reduce the cohomology of \(\mu_n\) to the groups \(\operatorname{Pic}(C)[n]\) and \(\operatorname{Pic}(C)/n\), both computed in Appendix \bpout{C}.

% object: proposition-smooth-proper-curves-torsion-prime (blueprint/chap-Cohomology-of-torsion-sheaves-on-curves/sec-Kummer-theory/proposition-smooth-proper-curves-torsion-prime/_section.md)
\phantomsection\label{bp10}%
\begin{bpstatementit}{Proposition~D.3}{ (Smooth proper curves, torsion prime to the characteristic)}
Let \(C\) be smooth, proper and connected over an algebraically closed field \(k\), with genus \(g\), and let \(n\geq1\) be invertible in \(k\). There are canonical identifications

\[
\begin{aligned}
H^0(C,\mu_n)&=\mu_n(k),&
H^1(C,\mu_n)&=J(C)[n](k),\\
H^2(C,\mu_n)&=\mathbf Z/n,&
H^q(C,\mu_n)&=0\quad(q>2).
\end{aligned}
\]

The isomorphism in degree two sends the Kummer boundary of a line bundle \(L\) to \(\deg L\bmod n\).  These identifications commute with extension of algebraically closed fields.
\end{bpstatementit}

% proof: proof/proposition-smooth-proper-curves-torsion-prime
\begin{proof}[{Proof}]\phantomsection\label{bp11}%
Pulling back the Kummer sequence of \(C_{\acute{e}t}\) by \(\nu_C^*\), we get the exact sequence of sheaves on \(C_{\mathrm{pro}\acute{e}t}\)

\[
0\longrightarrow\mu_n\longrightarrow\nu_C^*\mathbf G_m
\xrightarrow{(\cdot)^n}\nu_C^*\mathbf G_m
\longrightarrow0.
\]

Here surjectivity of the \(n\)th power map is a local question. For a unit \(a\) the equation \(T^n=a\) defines an étale cover, since its derivative \(nT^{n-1}\) is a unit. Moreover, the global units on the proper connected curve form the group \(k^\times\), which is \(n\)-divisible because \(k\) is algebraically closed. Hence the long exact sequence starts with \(H^0(C,\mu_n)=\mu_n(k)\). Since the map \(H^0(C,\mathbf G_m)\to H^0(C,\mathbf G_m)\) is onto, \(H^1(C,\mu_n)\) injects into \(H^1(C,\mathbf G_m)\). Using Lemma \hyperref[bp7]{D.2}, we read off from the long exact sequence that

\[
H^1(C,\mu_n)=\operatorname{Pic}(C)[n],\qquad
H^2(C,\mu_n)=\operatorname{Pic}(C)/n,
\]

with all higher groups zero.

Corollary \bpout{C.12} gives the exact sequence

\[
0\longrightarrow J(C)(k)\longrightarrow\operatorname{Pic}(C)
\xrightarrow{\deg}\mathbf Z\longrightarrow0.
\]

Multiplication by \(n\) on \(J(C)(k)\) is surjective, with kernel \(J(C)[n](k)\), and \(J(C)[n]\) is finite étale of rank \(n^{2g}\) (Corollary \bpout{C.15}).  Applying the snake lemma to multiplication by \(n\) on this sequence, and using that \(\mathbf Z\) is torsion free, we obtain \(\operatorname{Pic}(C)[n]=J(C)[n](k)\) and \(\operatorname{Pic}(C)/n=\mathbf Z/n\) via the degree. In particular, \(H^1(C,\mu_n)\) is noncanonically isomorphic to \((\mathbf Z/n)^{2g}\) [Stacks, \href{https://stacks.math.columbia.edu/tag/03RQ}{03RQ}].

Under an algebraically closed extension \(k'/k\), the schemes \(J(C)\) and \(J(C)[n]\) base change to the corresponding objects for \(C_{k'}\). The map \(J(C)[n](k)\to J(C)[n](k')\) is bijective by Corollary \bpout{C.15}. The same argument applies to the finite étale scheme \(\mu_n\). Finally, pullback of line bundles preserves the degree. Hence the pullback maps on cohomology are isomorphisms.
\end{proof}

% object: remark-twists-projective-line (blueprint/chap-Cohomology-of-torsion-sheaves-on-curves/sec-Kummer-theory/remark-twists-projective-line/_section.md)
\phantomsection\label{bp12}%
\begin{bpstatement}{Remark~D.4}{ (Twists and the projective line)}
The degree identifies \(H^2(C,\mu_n)\) with \(\mathbf Z/n\) without any choice. Since \(\mu_n=(\mathbf Z/n)(1)\), twisting by \(-1\) gives

\[
H^2(C,\mathbf Z/n)=(\mathbf Z/n)(-1),
\]

and identifying this group with \(\mathbf Z/n\) requires the choice of a primitive \(n\)th root of unity.  For \(C=\mathbf P^1_k\) the Jacobian is zero, so the cohomology with constant coefficients is concentrated in degrees zero and two, and the Kummer boundary of \(\mathcal O(1)\) corresponds to \(1\in\mathbf Z/n\). Section \bpout{4.4} uses this normalization on the geometric fibre.
\end{bpstatement}


% object: sec-Affine-curves (blueprint/chap-Cohomology-of-torsion-sheaves-on-curves/sec-Affine-curves/_section.md), content hidden
\subsection*{\texorpdfstring{D.2\quad Affine curves}{D.2 Affine curves}}
\phantomsection\label{bp13}%
\addcontentsline{toc}{subsection}{\texorpdfstring{D.2\quad Affine curves}{D.2 Affine curves}}
% text: prose-D-04
\phantomsection\label{bp14}%
Every smooth connected curve over an algebraically closed field that is not proper is affine and is the complement of a finite nonempty set of closed points in a smooth proper connected curve.  To see this, take a projective model of its function field and normalize it, and recall that regular curves over a perfect field are smooth.

On the affine curve \(U=\overline C\setminus D\) the Picard group is a quotient of \(\operatorname{Pic}^0(\overline C)\). It is therefore divisible, and \(H^2(U,\mu_n)=0\); this matches the fact that, topologically, a punctured Riemann surface is homotopy equivalent to a graph. In degree one, on the other hand, new torsors appear, since a Kummer torsor on \(U\) need not extend over \(\overline C\). The obstruction to extending it at each puncture is the residue of the torsor, that is, the order modulo \(n\) of the corresponding rational section. For example, on \(\mathbf G_m=\mathbf P^1\setminus\{0,\infty\}\) the torsor \(y^n=x\) has residues \(1\) and \(-1\) at the two punctures. It generates \(H^1(\mathbf G_m,\mu_n)\simeq\mathbf Z/n\). In general the residues sum to zero. The residue sequence of Proposition \hyperref[bp15]{D.5} shows that \(H^1(U,\mu_n)\) is free of rank \(2g+\#D-1\) over \(\mathbf Z/n\). This is the first Betti number of a surface of genus \(g\) with \(\#D\) punctures.

% object: proposition-cohomology-smooth-affine-curve (blueprint/chap-Cohomology-of-torsion-sheaves-on-curves/sec-Affine-curves/proposition-cohomology-smooth-affine-curve/_section.md)
\phantomsection\label{bp15}%
\begin{bpstatementit}{Proposition~D.5}{ (Cohomology of a smooth affine curve)}
Let \(\overline C\) be a smooth proper connected curve of genus \(g\) over an algebraically closed field \(k\).  Let \(U=\overline C\setminus D\), where \(D\) is a finite nonempty set of closed points, and let \(n\) be invertible in \(k\).  There is a natural exact sequence

\[
0\longrightarrow H^1(\overline C,\mu_n)
\longrightarrow H^1(U,\mu_n)
\xrightarrow{\partial}(\mathbf Z/n)^D
\xrightarrow{\sum}\mathbf Z/n\longrightarrow0.
\tag{D.1}
\]

Moreover, \(H^0(U,\mu_n)=\mu_n(k)\) and \(H^q(U,\mu_n)=0\) for \(q>1\). The group \(H^1(U,\mu_n)\) is noncanonically \((\mathbf Z/n)^{2g+\#D-1}\), and the pullback maps on all these groups are isomorphisms under algebraically closed extension.
\end{bpstatementit}

% proof: proof/proposition-cohomology-smooth-affine-curve
\begin{proof}[{Proof}]\phantomsection\label{bp16}%
By the Kummer sequence, the vanishing above degree one follows from the divisibility of \(\operatorname{Pic}(U)\) and the vanishing of \(H^q(U,\mathbf G_m)\) for \(q\geq2\) (Lemma \hyperref[bp7]{D.2}). For \(H^1\) we describe \(\mu_n\)-torsors as line bundles \(L\) with a trivialization of \(L^{\otimes n}\). After extending \(L\) to \(\overline C\), the orders at the punctures measure the failure of the torsor to extend. These orders give the residue map \(\partial\). We prove exactness at each place, and then deduce the rank and the invariance under field extension.

Since \(\overline C\) is regular, we may use divisors to see that

\[
\operatorname{Pic}(U)=\operatorname{Pic}(\overline C)/
\langle\mathcal O_{\overline C}(x):x\in D\rangle.
\]

One of these line bundles has degree one. Every class in \(\operatorname{Pic}(\overline C)\) is therefore congruent, modulo the subgroup generated by it, to a class of degree zero. Hence \(J(\overline C)(k)\) surjects onto \(\operatorname{Pic}(U)\). As \(J(\overline C)(k)\) is \(n\)-divisible (Corollary \bpout{C.15}), so is \(\operatorname{Pic}(U)\). By the Kummer sequence and Lemma \hyperref[bp7]{D.2}, we have \(H^2(U,\mu_n)=\operatorname{Pic}(U)/n=0\) and \(H^q(U,\mu_n)=0\) for \(q\geq3\).

We construct the sequence (D.1) using line bundles.  Recall that a \(\mu_n\)-torsor on \(U\) is given by a pair \((L,\alpha:L^{\otimes n}\simeq\mathcal O_U)\) of a line bundle and a trivialization of its \(n\)th power.  To define its residues, we extend \(L\) to a line bundle \(\overline L\) on \(\overline C\) and let \(s=\alpha^{-1}(1)\) be the corresponding rational section of \(\overline L^{\otimes n}\).  The divisor of \(s\) is supported on \(D\), since \(s\) generates on \(U\). Its orders modulo \(n\) do not depend on the extension. Indeed, two extensions differ by \(\mathcal O(E')\) with \(E'\) supported on \(D\), and this changes the divisor of \(s\) by \(nE'\). These orders define \(\partial(L,\alpha)\). Their sum is zero modulo \(n\), since the divisor of \(s\) has degree \(\deg\overline L^{\otimes n}=n\deg\overline L\).

Suppose all orders are divisible by \(n\), and write \(\operatorname{div}(s)=nE\) with \(E\) supported on \(D\). Then \(s\) is a nowhere vanishing section of \((\overline L(-E))^{\otimes n}=\overline L^{\otimes n}(-nE)\). So the trivialization extends to a torsor with line bundle \(\overline L(-E)\) on \(\overline C\), and this torsor restricts to \((L,\alpha)\) on \(U\). This proves exactness at \(H^1(U,\mu_n)\).

To prove injectivity, note that a trivialization on \(U\) of a torsor on \(\overline C\) gives a rational section of its line bundle whose \(n\)th power has no zero or pole on \(\overline C\).  Therefore the divisor of the section is zero, and the section trivializes the torsor on \(\overline C\).

To see that \(\partial\) maps onto the kernel of \(\sum\), we lift a residue vector with sum zero modulo \(n\) to an integral divisor \(B\) supported on \(D\). Its degree is divisible by \(n\) because the residues sum to zero. We claim that, as a consequence, there is a line bundle \(\overline L\) with \(\overline L^{\otimes n}\simeq\mathcal O_{\overline C}(B)\). To find it, choose \(P\in\overline C(k)\) and put \(L_1=\mathcal O_{\overline C}((\deg B/n)P)\), so that the line bundle \(\mathcal O_{\overline C}(B)\otimes L_1^{\otimes-n}\) has degree zero. By divisibility of \(J(\overline C)(k)\) it is \(N^{\otimes n}\) for some \(N\) of degree zero, and \(\overline L=L_1\otimes N\) works. It remains to produce the torsor. The canonical section of \(\mathcal O_{\overline C}(B)\) has divisor \(B\) and is nowhere zero on \(U\). It therefore defines a trivialization of \(\overline L^{\otimes n}|_U\), that is, a Kummer torsor on \(U\) whose residues are the orders of \(B\). Finally, the summation map is surjective since \(D\) is nonempty. This proves that the sequence is exact (see also [Stacks, \href{https://stacks.math.columbia.edu/tag/03RR}{03RR}]).

The left term and the kernel of summation are free \(\mathbf Z/n\)-modules of ranks \(2g\) and \(\#D-1\). Hence the resulting short exact sequence splits noncanonically. For naturality, note that the formation of the compactification, the boundary points, the line bundles and the orders of vanishing commutes with base change. The leftmost pullback map is an isomorphism by Proposition \hyperref[bp10]{D.3}, and the residue groups and their summation maps are unchanged. By the five lemma applied to the two exact sequences, the middle pullback map is an isomorphism.
\end{proof}

% text: prose-D-05
\phantomsection\label{bp17}%
For the constant sheaf \(\mathbf F_\ell\), with \(\ell\) a prime different from \(\operatorname{char}k\), the two preceding propositions prove \((C_1)\)–\((C_8)\) on smooth curves. Indeed, a primitive \(\ell\)th root of unity in \(k\) identifies \(\mathbf F_\ell\) with \(\mu_\ell\), compatibly with the pullback maps to \(C_{k'}\), since the root of unity stays primitive in \(k'\). For \((C_8)\), note that an open in a smooth connected proper curve is either the whole curve or an affine curve, which has \(H^2=0\), and that an open in a smooth affine curve is affine.


% object: sec-The-Artin-Schreier-sequence (blueprint/chap-Cohomology-of-torsion-sheaves-on-curves/sec-The-Artin-Schreier-sequence/_section.md), content hidden
\subsection*{\texorpdfstring{D.3\quad The Artin–Schreier sequence}{D.3 The Artin–Schreier sequence}}
\phantomsection\label{bp18}%
\addcontentsline{toc}{subsection}{\texorpdfstring{D.3\quad The Artin–Schreier sequence}{D.3 The Artin–Schreier sequence}}
% text: prose-D-06
\phantomsection\label{bp19}%
Since the Kummer cover \(T^n=a\) is étale only when \(n\) is invertible, \(p\)-torsion in characteristic \(p\) needs a different sequence. In the Artin–Schreier sequence the role of \(\mathbf G_m\) is played by the additive group, whose étale cohomology is the coherent cohomology of the structure sheaf. Since coherent cohomology of a curve vanishes above degree one, this gives \(H^q(C,\mathbf F_p)=0\) for \(q\geq3\). For \(q=2\) we also need that \(\varphi-1\) is surjective on \(H^1(C,\mathcal O_C)\), which holds over an algebraically closed field. So \(H^2(C,\mathbf F_p)=0\) even on a proper curve, and there is no analogue of the class of a point in \(H^2\).

On a proper curve, \(H^1(C,\mathbf F_p)\) is the fixed space of the Frobenius-linear map \(\varphi\) on the \(g\)-dimensional space \(H^1(C,\mathcal O_C)\). Hence it has dimension at most \(g\) (Proposition \hyperref[bp20]{D.6}). The bound can be strict. For example, for an elliptic curve, \(\varphi\) on the line \(H^1(E,\mathcal O_E)\) is bijective when \(E\) is ordinary, and zero when \(E\) is supersingular, as is \(y^2=x^3-x\) over \(\overline{\mathbf F}_3\). Accordingly \(H^1(E,\mathbf F_p)\) has dimension one or zero. On an affine curve, by contrast, the cokernel of \(\varphi-1\) on the ring of functions contributes to \(H^1\), and this group is infinite (Example \hyperref[bp26]{D.8}).

% object: proposition-p-torsion-smooth-proper-curve (blueprint/chap-Cohomology-of-torsion-sheaves-on-curves/sec-The-Artin-Schreier-sequence/proposition-p-torsion-smooth-proper-curve/_section.md)
\phantomsection\label{bp20}%
\begin{bpstatementit}{Proposition~D.6}{ ($p$-torsion on a smooth proper curve in characteristic $p$)}
Let \(C\) be a smooth proper connected curve of genus \(g\) over an algebraically closed field of characteristic \(p>0\), and let \(\varphi\) be the Frobenius-linear endomorphism of \(H^1(C,\mathcal O_C)\) induced by \(a\mapsto a^p\) on \(\mathcal O_C\).  Then

\[
H^0(C,\mathbf F_p)=\mathbf F_p,\qquad
H^1(C,\mathbf F_p)=
\ker(\varphi-1:H^1(C,\mathcal O_C)\to H^1(C,\mathcal O_C)),
\]

and \(H^q(C,\mathbf F_p)=0\) for \(q>1\).  The group \(H^1(C,\mathbf F_p)\) is finite, of \(\mathbf F_p\)-dimension at most \(g\), and its pullback map is an isomorphism under algebraically closed extension.
\end{bpstatementit}

% proof: proof/proposition-p-torsion-smooth-proper-curve
\begin{proof}[{Proof}]\phantomsection\label{bp21}%
The Artin–Schreier sequence turns the question into one about \(\varphi-1\) on \(H^0(C,\mathcal O_C)=k\) and on \(H^1(C,\mathcal O_C)\). The key fact is that on a finite-dimensional space over an algebraically closed field, a Frobenius-linear map minus the identity is surjective with finite kernel.

\emph{Step 1: the Artin–Schreier sequence.} Write also \(\varphi\) for the endomorphism \(a\mapsto a^p\) of the sheaf \(\mathbf G_a\) on \(C_{\acute{e}t}\).  The Artin–Schreier sequence

\[
0\longrightarrow\mathbf F_p\longrightarrow\mathbf G_a
\xrightarrow{\varphi-1}\mathbf G_a\longrightarrow0
\]

is exact, because surjectivity is a local question and \(T^p-T-a\) has derivative \(-1\). The étale cohomology of \(\mathbf G_a\) is the coherent cohomology of \(\mathcal O_C\). Indeed, quasi-coherent sheaves have the same cohomology on the étale and on the Zariski site, by faithfully flat descent on affine opens and the Čech resolution. This coherent cohomology vanishes in degrees above one, since \(C\) is a curve. Since \(C\) is proper, \(H^1(C,\mathcal O_C)\) is a finite-dimensional \(k\)-vector space.

\emph{Step 2: linear algebra.} Let \(\varphi\) be a Frobenius-linear endomorphism of a finite-dimensional \(k\)-vector space \(V\); we show that \(\varphi-1\) is surjective and that its kernel is the set of \(k\)-points of a finite étale scheme of order \(p^f\) with \(f\leq\dim_k V\).  To prove this, we reduce along a finite filtration to the part where \(\varphi\) is nilpotent and the part where \(\varphi\) is bijective.  On the nilpotent part, \(\varphi-1\) is an automorphism given by an additive polynomial, with inverse \(-(1+\varphi+\varphi^2+\cdots)\), a finite sum; in particular its kernel there is zero.  On a bijective part of dimension \(f\), by contrast, we choose a basis, invert the matrix of \(\varphi\), and write the equations of the fibres as

\[
x_i^p-\sum_j b_{ij}x_j=c_i\qquad(1\leq i\leq f),
\]

where \((b_{ij})\) is invertible.  The coordinate algebra is free of dimension \(p^f\), with basis the monomials having each exponent less than \(p\).  Moreover, the Jacobian matrix of these equations is \(-(b_{ij})\), because \(p=0\) in \(k\), and it is invertible; so each fibre is finite étale of degree \(p^f\) and has \(p^f\) points over the algebraically closed field \(k\) [Stacks, \href{https://stacks.math.columbia.edu/tag/0A3L}{0A3L}]. In particular every fibre is nonempty, so \(\varphi-1\) is surjective, and the kernel has \(p^f\) elements.

\emph{Step 3: conclusion.} We apply Step 2 to \(H^0(C,\mathcal O_C)=k\) and to \(H^1(C,\mathcal O_C)\) in the long exact sequence associated to the Artin–Schreier sequence. Surjectivity of \(\varphi-1\) on \(H^0\) gives \(H^1(C,\mathbf F_p)=\ker(\varphi-1)\) on \(H^1(C,\mathcal O_C)\), with the bound on its dimension, and surjectivity on \(H^1\) together with \(H^2(C,\mathcal O_C)=0\) gives \(H^q(C,\mathbf F_p)=0\) for \(q>1\). Finally, the pullback map under an extension \(k'/k\) is an isomorphism [Stacks, \href{https://stacks.math.columbia.edu/tag/0HAN}{0HAN}, \href{https://stacks.math.columbia.edu/tag/0A3P}{0A3P}]: flat base change for coherent cohomology identifies the vector space and its Frobenius-linear map with their scalar extensions, and the kernel is a finite étale scheme, whose \(k\)-points and \(k'\)-points agree. As always, Theorem \bpout{2.35} carries the statement over to \(C_{\mathrm{pro}\acute{e}t}\).
\end{proof}

% text: prose-D-07
\phantomsection\label{bp22}%
On an affine curve the same sequence is even simpler to use, because coherent cohomology vanishes in positive degrees. On the other hand, \(H^1\) is now the cokernel of \(\varphi-1\) on an infinite-dimensional space.

% object: lemma-p-torsion-affine-curve-characteristic (blueprint/chap-Cohomology-of-torsion-sheaves-on-curves/sec-The-Artin-Schreier-sequence/lemma-p-torsion-affine-curve-characteristic/_section.md)
\phantomsection\label{bp23}%
\begin{bpstatementit}{Lemma~D.7}{ ($p$-torsion on an affine curve in characteristic $p$)}
If \(U\) is affine over an algebraically closed field \(k\) of characteristic \(p\), then, for the classical constant sheaf \(\mathbf F_p\) and \(\varphi(a)=a^p\) on \(\Gamma(U,\mathcal O_U)\), we have

\[
H^1(U,\mathbf F_p)=
\Gamma(U,\mathcal O_U)/(\varphi-1)\Gamma(U,\mathcal O_U),\qquad
H^q(U,\mathbf F_p)=0\quad(q>1).
\]
\end{bpstatementit}

% proof: proof/lemma-p-torsion-affine-curve-characteristic
\begin{proof}[{Proof}]\phantomsection\label{bp24}%
Quasi-coherent cohomology on an affine scheme vanishes in positive degrees, and, as in the proof of the proposition, étale and Zariski cohomology of \(\mathbf G_a\) agree. In the long exact sequence associated to the Artin–Schreier sequence on \(U_{\acute{e}t}\), the terms \(H^q(U,\mathbf G_a)\) for \(q\geq1\) are therefore zero. What remains is the exact sequence \(0\to H^0(U,\mathbf F_p)\to\Gamma(U,\mathcal O_U)\xrightarrow{\varphi-1}\Gamma(U,\mathcal O_U)\to H^1(U,\mathbf F_p)\to0\), together with \(H^q(U,\mathbf F_p)=0\) for \(q>1\).  By Theorem \bpout{2.35} these statements carry over to \(U_{\mathrm{pro}\acute{e}t}\).
\end{proof}

% text: prose-D-08
\phantomsection\label{bp25}%
For the affine line we can compute the cokernel explicitly, and we see that it is infinite and grows with the field, unlike in the proper case.

% object: example-affine-line-characteristic-p (blueprint/chap-Cohomology-of-torsion-sheaves-on-curves/sec-The-Artin-Schreier-sequence/example-affine-line-characteristic-p/_section.md)
\phantomsection\label{bp26}%
\begin{bpstatement}{Example~D.8}{ (The affine line in characteristic $p$)}
For \(U=\mathbf A^1_k\) in characteristic \(p\) we get

\[
H^1(U,\mathbf F_p)=k[t]/\{h^p-h:h\in k[t]\}.
\]

Every class has a unique representative with zero constant term and all positive exponents prime to \(p\).  Indeed, we can remove exponents divisible by \(p\) using \(p\)th roots in the perfect field, and we can remove the constant term since \(a\mapsto a^p-a\) is surjective on \(k\). For uniqueness, note that a nonconstant polynomial \(h^p-h\) has degree divisible by \(p\).  In particular the quotient is infinite, because the classes of \(t^m\) for positive \(m\) prime to \(p\) are \(\mathbf F_p\)-linearly independent.  For an algebraically closed extension \(k\subsetneq k'\) and \(a\in k'\setminus k\), the class of \(at\) is not in the image of the quotient over \(k\). Indeed, the reduction of a polynomial in \(k[t]\) uses only \(p\)th roots and roots of \(x^p-x=c\) in \(k\), so its reduced representative lies in \(k[t]\), whereas \(at\) is already reduced. Thus finiteness and invariance under field extension in degree one fail for \(p\)-torsion coefficients on an affine curve.
\end{bpstatement}

% text: prose-D-09
\phantomsection\label{bp27}%
For the constant sheaf \(\mathbf F_p\) on smooth curves, we obtain \((C_1)\), \((C_3)\), \((C_5)\) and \((C_7)\) from the proposition on proper curves, and \((C_1)\), \((C_2)\) and \((C_3)\) from the lemma on affine curves, while \((C_8)\) holds because \(H^2\) vanishes in both cases.


% object: sec-Normalization (blueprint/chap-Cohomology-of-torsion-sheaves-on-curves/sec-Normalization/_section.md), content hidden
\subsection*{\texorpdfstring{D.4\quad Normalization}{D.4 Normalization}}
\phantomsection\label{bp28}%
\addcontentsline{toc}{subsection}{\texorpdfstring{D.4\quad Normalization}{D.4 Normalization}}
% text: prose-D-10
\phantomsection\label{bp29}%
So far we have computed with constant coefficients on smooth curves, whereas Theorem \bpout{4.20} concerns arbitrary curves of dimension at most one and arbitrary torsion sheaves. Two reductions close the gap: normalization and removing finitely many points. Both rest on Lemma \hyperref[bp33]{D.10}, which says that a sheaf supported on finitely many points has no higher cohomology. Hence passing between a curve, its normalization and a dense open subset changes cohomology only in degrees zero and one, and only by finite groups that do not change under field extension.

For example, let \(C:y^2=x^3+x^2\) be the nodal cubic. Its normalization \(\mathbf A^1\to C\) has two points over the node. For \(M=\mathbf Z/n\) the normalization sequence of Lemma \hyperref[bp36]{D.11} has cokernel \(M^2/\Delta M\simeq M\) at the node. If \(n\) is invertible in \(k\), then \(H^1(\mathbf A^1,\mathbf Z/n)=0\), and the long exact sequence gives \(H^1(C,\mathbf Z/n)\simeq\mathbf Z/n\). This is the analogue of the loop created by the node.

% object: lemma-finite-morphisms-cohomology (blueprint/chap-Cohomology-of-torsion-sheaves-on-curves/sec-Normalization/lemma-finite-morphisms-cohomology/_section.md)
\phantomsection\label{bp30}%
\begin{bpstatementit}{Lemma~D.9}{ (Finite morphisms and cohomology)}
Let \(h:Y\to C\) be a finite morphism of schemes. The direct image \(h_*\) of classical sheaves is exact, and \(H^q(C,h_*G)=H^q(Y,G)\) for every classical abelian sheaf \(G\) on \(Y_{\mathrm{pro}\acute{e}t}\) and every \(q\).
\end{bpstatementit}

% proof: proof/lemma-finite-morphisms-cohomology
\begin{proof}[{Proof}]\phantomsection\label{bp31}%
On the étale sites, Lemma \hyperref[bp3]{D.1} shows that \(h_*\) is exact and that its stalks are finite products of stalks of \(G\).  Since \(h_*\) is exact, the Leray spectral sequence gives \(H^q(C,h_*G)=H^q(Y,G)\).  The comparison of classical direct images and Theorem \bpout{2.35} transport these statements to the pro-étale site; see [Stacks, \href{https://stacks.math.columbia.edu/tag/03QP}{03QP}, \href{https://stacks.math.columbia.edu/tag/0959}{0959}].
\end{proof}

% text: prose-D-11
\phantomsection\label{bp32}%
Since a sheaf supported on finitely many points is a direct image from those points, it has no higher cohomology. All the reductions below use this vanishing.

% object: lemma-sheaves-finite-support-finite-morphisms (blueprint/chap-Cohomology-of-torsion-sheaves-on-curves/sec-Normalization/lemma-sheaves-finite-support-finite-morphisms/_section.md)
\phantomsection\label{bp33}%
\begin{bpstatementit}{Lemma~D.10}{ (Sheaves with finite support)}
Let \(F\) be a classical abelian sheaf supported on finitely many closed points of a scheme \(C\) of finite type over an algebraically closed field \(k\). Then \(H^q(C,F)=0\) for \(q>0\), and \(H^0(C,F)\) is the direct sum of the stalks of \(F\) at these points; the pullback map on \(H^0\) is an isomorphism under algebraically closed extension of \(k\).
\end{bpstatementit}

% proof: proof/lemma-sheaves-finite-support-finite-morphisms
\begin{proof}[{Proof}]\phantomsection\label{bp34}%
The sheaf \(F\) is \(h_*\) of a sheaf on a finite disjoint union of copies of \(\operatorname{Spec}k\), where \(h\) is the inclusion of the points. On these points taking sections is exact because \(k\) is algebraically closed, and the sections are the product, here a finite direct sum, of the stalks, which do not change under extension of \(k\). Finally, Lemma \hyperref[bp30]{D.9} carries these assertions over to \(C\).
\end{proof}

% text: prose-D-12
\phantomsection\label{bp35}%
The normalization of a curve is finite and is an isomorphism away from finitely many points, so the constant sheaf on \(C\) differs from the direct image of the constant sheaf on \(C^{\mathrm{nor}}\) only by a sheaf with finite support.

% object: lemma-normalization-constant-sheaf (blueprint/chap-Cohomology-of-torsion-sheaves-on-curves/sec-Normalization/lemma-normalization-constant-sheaf/_section.md)
\phantomsection\label{bp36}%
\begin{bpstatementit}{Lemma~D.11}{ (Normalization of a constant sheaf)}
Let \(C\) be a reduced curve of finite type over an algebraically closed field, and let \(\pi:C^{\mathrm{nor}}\to C\) be its normalization.  For a finite abelian group \(M\), the natural map

\[
\underline M_C\longrightarrow \pi_*\underline M_{C^{\mathrm{nor}}}
\]

is injective with cokernel supported on finitely many points.  At a point with \(b\) geometric branches, the cokernel is \(M^b/\Delta M\), where \(\Delta\) is the diagonal. These cokernels and the resulting exact sequences commute with algebraically closed field extension.
\end{bpstatementit}

% proof: proof/lemma-normalization-constant-sheaf
\begin{proof}[{Proof}]\phantomsection\label{bp37}%
The map is computed on stalks by the formula for stalks of finite direct images, and the cokernel has finite support because \(\pi\) is an isomorphism away from finitely many points. Compatibility with field extension holds because normalization commutes with algebraically closed extension.

Normalization is finite for reduced schemes of finite type over a field.  On each one-dimensional component, its normal model is regular, hence smooth over the perfect field, while the isolated components of dimension zero are reduced points, that is, copies of the spectrum of the ground field.  Since the normalization is an isomorphism over a dense open of each component, the locus where it is not an isomorphism is finite.

By the calculation of stalks in Lemma \hyperref[bp3]{D.1}, the displayed map at a geometric point is the diagonal \(M\to M^b\), with \(b\geq1\), so its kernel is zero and its cokernel has the asserted form.  Where \(\pi\) is an isomorphism we have \(b=1\) and the cokernel vanishes; hence the cokernel is supported on the finite locus found above.

After algebraically closed extension, the normalization remains smooth and is a finite birational normal model of the reduced base change, so it is the normalization of the base change.  Base change for finite morphisms identifies the fibres and their sets of branches, and therefore the diagonal maps and their cokernels are obtained by extension of scalars.
\end{proof}

% text: prose-D-13
\phantomsection\label{bp38}%
The normalization sequence reduces constant coefficients on a reduced curve to the smooth case. By Lemma \hyperref[bp33]{D.10}, its cokernel term has finite sections, unchanged under field extension, and no positive cohomology. Moreover, by Lemma \hyperref[bp30]{D.9}, \(H^q(C,\pi_*\underline M)=H^q(C^{\mathrm{nor}},\underline M)\). The long exact sequence therefore yields isomorphisms \(H^q(C,\underline M)\simeq H^q(C^{\mathrm{nor}},\underline M)\) for \(q\geq2\), while in degrees zero and one it gives a natural exact sequence relating \(C\) and \(C^{\mathrm{nor}}\), whose remaining term is a finite group. From these sequences, vanishing and finiteness pass from \(C^{\mathrm{nor}}\) to \(C\), and by the five lemma so does invariance under field extension.

The normalization of an affine or proper curve is again affine or proper. Hence all assertions of Theorem \bpout{4.20} for constant \(\mathbf F_\ell\) follow from the calculations on smooth curves. For \((C_8)\), we apply this comparison to \(C\) and to the inverse image of an open subset. The commutative square of restriction maps then reduces surjectivity to the smooth components.


% object: sec-Extension-by-zero (blueprint/chap-Cohomology-of-torsion-sheaves-on-curves/sec-Extension-by-zero/_section.md), content hidden
\subsection*{\texorpdfstring{D.5\quad Extension by zero}{D.5 Extension by zero}}
\phantomsection\label{bp39}%
\addcontentsline{toc}{subsection}{\texorpdfstring{D.5\quad Extension by zero}{D.5 Extension by zero}}
% text: prose-D-14
\phantomsection\label{bp40}%
A constructible sheaf on a curve is locally constant on a dense open subset. Removing the finitely many remaining points replaces the sheaf by its extension by zero from this open, and changes cohomology only in degrees zero and one. Let \(j:U\hookrightarrow C\) be a dense open in a reduced curve, and let \(i:Z\hookrightarrow C\) be its finite complementary closed subscheme. For a classical sheaf \(F\) we have the exact sequence

\[
0\longrightarrow j_!j^*F\longrightarrow F
\longrightarrow i_*i^*F\longrightarrow0
\]

of an open and its closed complement. We check its exactness on geometric stalks on \(C_{\acute{e}t}\), and then pull it back.

% object: lemma-removing-finitely-many-points (blueprint/chap-Cohomology-of-torsion-sheaves-on-curves/sec-Extension-by-zero/lemma-removing-finitely-many-points/_section.md)
\phantomsection\label{bp41}%
\begin{bpstatementit}{Lemma~D.12}{ (Removing finitely many points)}
Let \(C\) be a reduced separated curve of finite type over an algebraically closed field \(k\), let \(j:U\hookrightarrow C\) be a dense open with finite complement \(i:Z\hookrightarrow C\), let \(\ell\) be a prime, and let \(F\) be a classical constructible sheaf with \(\ell F=0\). For each of the assertions \((C_1)\)–\((C_8)\) listed at the beginning of this appendix, under its hypotheses on \(C\) and on the prime \(\ell\), the assertion holds for \(F\) if and only if it holds for \(j_!j^*F\).
\end{bpstatementit}

% proof: proof/lemma-removing-finitely-many-points
\begin{proof}[{Proof}]\phantomsection\label{bp42}%
The third term of the exact sequence displayed above is supported on the finite set \(Z\), so by Lemma \hyperref[bp33]{D.10} it has finite sections and no higher cohomology. Comparing the long exact sequences of \(F\) and \(j_!j^*F\) then shows that each assertion passes in both directions.

The sheaf \(j_!j^*F\) is again killed by \(\ell\), so the hypotheses on the coefficient prime are the same for both sheaves. Since the last term of the exact sequence has no positive cohomology, the long exact sequence yields

\[
\begin{split}
0\to H^0(C,j_!j^*F)\to H^0(C,F)\to H^0(Z,i^*F)\to{}\qquad\qquad\\
H^1(C,j_!j^*F)\to H^1(C,F)\to0,
\end{split}
\]

and isomorphisms in degrees at least two. Since the sequence is natural under field extension and its term supported on \(Z\) is unchanged, the two long exact sequences show that the pullback maps are isomorphisms for \(F\) in all degrees if and only if they are isomorphisms for \(j_!j^*F\) in all degrees. For \((C_8)\), we apply the degree-two identifications after restriction to every open subset of \(C\). See also [Stacks, \href{https://stacks.math.columbia.edu/tag/0GJA}{0GJA}].
\end{proof}

% text: prose-D-15
\phantomsection\label{bp43}%
When \(j\) is dense, Lemma \hyperref[bp41]{D.12}, together with the calculation for constant \(\mathbf F_\ell\) on \(C\) via Lemma \hyperref[bp36]{D.11}, shows that \(j_!\underline{\mathbf F_\ell}\) satisfies the assertions of Theorem \bpout{4.20}. An arbitrary open is dense in the reduced union of the components it meets, a closed subscheme of \(C\), so we can apply the dense case there and pass back to \(C\): direct image along a finite morphism is exact [Stacks, \href{https://stacks.math.columbia.edu/tag/0A3Q}{0A3Q}] and, by Lemma \hyperref[bp30]{D.9}, does not change cohomology.


% object: sec-Reduction-to-constant-coefficients (blueprint/chap-Cohomology-of-torsion-sheaves-on-curves/sec-Reduction-to-constant-coefficients/_section.md), content hidden
\subsection*{\texorpdfstring{D.6\quad Reduction to constant coefficients}{D.6 Reduction to constant coefficients}}
\phantomsection\label{bp44}%
\addcontentsline{toc}{subsection}{\texorpdfstring{D.6\quad Reduction to constant coefficients}{D.6 Reduction to constant coefficients}}
% text: prose-D-16
\phantomsection\label{bp45}%
We reduce a local system to constant sheaves by pulling it back along a cover of degree prime to \(\ell\), on which it becomes unipotent. To find such a cover, note that a local system of finite \(\mathbf F_\ell\)-vector spaces has finite monodromy: the fundamental group acts on a fibre through a finite quotient. The local system then becomes a successive extension of constant sheaves on the cover defined by an \(\ell\)-Sylow subgroup of the monodromy group.

On the full cover of its monodromy the local system even becomes constant. But that cover can have degree divisible by \(\ell\), and then the local system need not be a direct summand of the pushforward of its pullback. For example, let \(v:y^2=x\) be the double cover of \(\mathbf G_m\), in characteristic different from two, and let \(\ell=2\). The composite of the unit \(\mathbf F_2\to v_*\mathbf F_2\) and the trace is multiplication by \(2\), which is zero. In fact \(\mathbf F_2\) is not a direct summand of \(v_*\mathbf F_2\) at all. The monodromy representation of \(v_*\mathbf F_2\) is the regular representation \(\mathbf F_2[\mathbf Z/2]\), which is indecomposable.

For a cover of degree prime to \(\ell\), on the other hand, Lemma \hyperref[bp46]{D.13} shows that the trace divided by the degree splits the unit. It therefore suffices to prove the vanishing, finiteness and invariance statements for the pullback. The pullback is unipotent, that is, built from constant sheaves by extensions.

% object: lemma-reduction-unipotent-local-system (blueprint/chap-Cohomology-of-torsion-sheaves-on-curves/sec-Reduction-to-constant-coefficients/lemma-reduction-unipotent-local-system/_section.md)
\phantomsection\label{bp46}%
\begin{bpstatementit}{Lemma~D.13}{ (Reduction to a unipotent local system)}
Let \(\ell\) be a prime, let \(U\) be a connected scheme, and let \(G\) be a finite classical local system of \(\mathbf F_\ell\)-vector spaces.  There is a finite étale cover \(v:V\to U\) of degree \(d\) prime to \(\ell\) such that \(v^*G\) has a finite filtration with successive quotients the constant sheaf \(\mathbf F_\ell\).  Moreover, \(G\) is a direct summand of \(v_*v^*G\).
\end{bpstatementit}

% proof: proof/lemma-reduction-unipotent-local-system
\begin{proof}[{Proof}]\phantomsection\label{bp47}%
The filtration comes from the representation theory of \(\ell\)-groups, whose only simple module over \(\mathbf F_\ell\) is the trivial one, and the splitting from the trace, whose composite with the unit is multiplication by the degree.

By Proposition \bpout{2.42}, \(G\) is the pullback of a finite étale local system, which corresponds to a finite continuous representation of the fundamental group.  Let \(H\) be the finite image of this representation, choose an \(\ell\)-Sylow subgroup \(P\), and take the cover associated to its inverse image in the fundamental group.  Its degree is \([H:P]\), which is prime to \(\ell\), and the restricted action factors through \(P\).  Since the augmentation ideal of \(\mathbf F_\ell[P]\) is nilpotent, the simple modules of the finite \(\ell\)-group \(P\) are trivial of dimension one, and a composition series yields the required filtration [Stacks, \href{https://stacks.math.columbia.edu/tag/0A3R}{0A3R}].

On a fibre, the unit \(G\to v_*v^*G\) is the diagonal and the trace map \(\operatorname{Tr}:v_*v^*G\to G\) adds the \(d\) entries. Their composite is therefore multiplication by \(d\), which is invertible on an \(\mathbf F_\ell\)-sheaf, so \(d^{-1}\operatorname{Tr}\) is a retraction.
\end{proof}

% text: prose-D-17
\phantomsection\label{bp48}%
To apply this to \(j_!G\) on a reduced curve \(C\), we take for \(U\) a dense smooth open on which \(G\) is locally constant.  Normalizing \(C\) in the finite extensions of function fields defined by \(V\), we get a finite morphism \(\bar v:Y\to C\) with \(V\) the inverse image of \(U\), and an open immersion \(j':V\hookrightarrow Y\).  From the cartesian square formed by these maps and the formula for stalks of direct images along finite morphisms, we find that

\[
j_!v_*v^*G=\bar v_*j'_!v^*G,
\]

which is also the exchange formula of [Stacks, \href{https://stacks.math.columbia.edu/tag/0GJ1}{0GJ1}].  If \(C\) is proper, then \(Y\) is proper, and if \(C\) is affine, then \(Y\) is affine.

The assertions travel from constant sheaves to \(j'_!v^*G\) on \(Y\), then to \(j_!v_*v^*G\) on \(C\), and from there to the direct summand \(j_!G\). Since \(j'_!\) is exact, the filtration of \(v^*G\) gives a filtration of \(j'_!v^*G\) with graded pieces \(j'_!\mathbf F_\ell\). The assertions of Theorem \bpout{4.20} hold for these pieces by the calculation for constant sheaves extended by zero. As we show at the start of the next section, they pass along extensions, so they hold for \(j'_!v^*G\). Since \(\bar v_*\) is exact, they hold for \(j_!v_*v^*G\). By the splitting given by the trace, they hold for its direct summand \(j_!G\). For invariance under field extension we use that the unit and the trace commute with extension of the ground field. Then the pullback map for \(j_!G\) is a retract of the pullback map for \(j_!v_*v^*G\), and a retract of an isomorphism is an isomorphism [Stacks, \href{https://stacks.math.columbia.edu/tag/03SD}{03SD}].

If \(C\) is not equidimensional, we choose a dense open that is smooth and is the disjoint union of the regular loci of the one-dimensional components and the isolated points.  On the isolated points taking sections is exact, each connected one-dimensional piece is treated as above, and finite direct sums preserve all the assertions.


% object: sec-Devissage (blueprint/chap-Cohomology-of-torsion-sheaves-on-curves/sec-Devissage/_section.md), content hidden
\subsection*{\texorpdfstring{D.7\quad Dévissage}{D.7 Dévissage}}
\phantomsection\label{bp49}%
\addcontentsline{toc}{subsection}{\texorpdfstring{D.7\quad Dévissage}{D.7 Dévissage}}
% text: prose-D-18
\phantomsection\label{bp50}%
The assertions of Theorem \bpout{4.20} pass along short exact sequences.  To see this, let \(0\to F_1\to F\to F_2\to0\) be exact, and suppose that the assertions are known for \(F_1\) and \(F_2\).  Vanishing and finiteness follow from the long exact sequence, and invariance under field extension follows from the map between the two long exact sequences and the five lemma. For \((C_8)\) an extra argument is needed. Let \(b\in H^2(U,F)\), and lift its image in \(H^2(U,F_2)\) to \(c_2\in H^2(C,F_2)\).  The obstruction to lifting \(c_2\) to \(H^2(C,F)\) lies in \(H^3(C,F_1)\), which vanishes by \((C_1)\), so such a lift \(c\) exists. The difference \(b-c|_U\) comes from \(H^2(U,F_1)\), and by \((C_8)\) for \(F_1\) it lifts to \(H^2(C,F_1)\); adding that lift to \(c\) gives the required global class [Stacks, \href{https://stacks.math.columbia.edu/tag/0A5C}{0A5C}].

% proof: proof/theorem-cohomology-curves/2
\begin{proof}[{Proof of Theorem 4.20}]\phantomsection\label{bp51}%
For constructible sheaves we reduce to the explicit calculations, and arbitrary torsion sheaves are filtered colimits of these. The reduction has four steps: dévissage to constructible \(\mathbf F_\ell\)-sheaves, extension by zero from an open on which the sheaf is locally constant, a cover of degree prime to \(\ell\) that makes it unipotent, and normalization, which leaves constant coefficients on smooth curves.

By Theorem \bpout{2.35}, we may work on \(C_{\acute{e}t}\), and since a universal homeomorphism does not change the étale topos, we may replace \(C\) by \(C_{\mathrm{red}}\).  A constructible torsion sheaf has a finite exponent, because it has finitely many isomorphism types of finite stalks.  It is therefore the direct sum of its finitely many primary parts, and a finite filtration of an \(\ell\)-primary part has constructible \(\mathbf F_\ell\)-sheaves as graded pieces.  By the stability under extensions just proved, it suffices to treat such sheaves.

Choose a dense smooth open \(j:U\hookrightarrow C\) on which the sheaf is finite locally constant.  By Lemma \hyperref[bp41]{D.12} we may replace the sheaf by its extension by zero.  On the connected components of \(U\), Lemma \hyperref[bp46]{D.13} and finite direct images reduce us to constant \(\mathbf F_\ell\) extended by zero on curves finite over \(C\); this is allowed because the assertions pass to direct summands and because, by Lemma \hyperref[bp30]{D.9}, finite direct images are exact and do not change cohomology.  Finally, normalization and Lemma \hyperref[bp33]{D.10} reduce these sheaves to the calculations with constant coefficients on smooth curves, namely the Kummer sequence for \(\ell\ne p\) and the Artin–Schreier sequence for \(\ell=p\).  Since all reductions preserve the hypothesis that \(C\) is proper or affine, the theorem holds for constructible torsion sheaves.

An arbitrary torsion sheaf on the qcqs scheme \(C\) is a filtered colimit of constructible torsion sheaves, and cohomology commutes with filtered colimits [Stacks, \href{https://stacks.math.columbia.edu/tag/03SA}{03SA}, \href{https://stacks.math.columbia.edu/tag/03Q5}{03Q5}].  For torsion prime to the characteristic or \(p\)-primary torsion, we may take the constructible sheaves with the same kind of torsion, by splitting off primary parts before passing to the colimit. Pullback and restriction also commute with filtered colimits. Hence vanishing and the pullback isomorphisms pass to the colimit. So does surjectivity in \((C_8)\), because a representative of a class already occurs at some stage. Finiteness is asserted only for constructible sheaves. Finally, Theorem \bpout{2.35} carries these conclusions and their maps over to \(C_{\mathrm{pro}\acute{e}t}\).
\end{proof}



\clearpage\phantomsection\hypertarget{bp-outside}{}%
\section*{Outside this export}
\addcontentsline{toc}{section}{Outside this export}
References marked \ensuremath{{}^\ast} are to parts of the blueprint this document leaves out:
\begin{description}
\item[Definition 2.30] Classical pro-étale sheaves, in Ⅰ The six operations (\bpslug{definition-{}\allowbreak{}classical-{}\allowbreak{}pro-{}\allowbreak{}etale-{}\allowbreak{}sheaves})
\item[Theorem 2.32] Classical sheaves as the essential image, in Ⅰ The six operations (\bpslug{theorem-{}\allowbreak{}classical-{}\allowbreak{}sheaves-{}\allowbreak{}essential-{}\allowbreak{}image})
\item[Theorem 2.35] Bounded-below cohomology comparison, in Ⅰ The six operations (\bpslug{theorem-{}\allowbreak{}bounded-{}\allowbreak{}below-{}\allowbreak{}cohomology-{}\allowbreak{}comparison})
\item[Proposition 2.42] Locally constant sheaves of finite presentation, in Ⅰ The six operations (\bpslug{proposition-{}\allowbreak{}locally-{}\allowbreak{}constant-{}\allowbreak{}sheaves-{}\allowbreak{}finite-{}\allowbreak{}presentation})
\item[Section 4.4] The projective line, in Ⅰ The six operations (\bpslug{sec-{}\allowbreak{}The-{}\allowbreak{}projective-{}\allowbreak{}line})
\item[Theorem 4.20] Cohomology of curves, in Ⅰ The six operations (\bpslug{theorem-{}\allowbreak{}cohomology-{}\allowbreak{}curves})
\item[Lemma 4.21] Separably closed field extension, in Ⅰ The six operations (\bpslug{lemma-{}\allowbreak{}extension-{}\allowbreak{}separably-{}\allowbreak{}closed-{}\allowbreak{}ground-{}\allowbreak{}field})
\item[Appendix C] The Picard scheme of a curve, in Appendices (\bpslug{chap-{}\allowbreak{}The-{}\allowbreak{}Picard-{}\allowbreak{}scheme-{}\allowbreak{}of-{}\allowbreak{}a-{}\allowbreak{}curve})
\item[Theorem C.3] Riemann–Roch and coherent Serre duality, in Appendices (\bpslug{theorem-{}\allowbreak{}riemann-{}\allowbreak{}roch-{}\allowbreak{}coherent-{}\allowbreak{}serre-{}\allowbreak{}duality})
\item[Corollary C.12] Degree and the Jacobian, in Appendices (\bpslug{corollary-{}\allowbreak{}degree-{}\allowbreak{}jacobian})
\item[Corollary C.15] Torsion of the Jacobian and field extension, in Appendices (\bpslug{corollary-{}\allowbreak{}torsion-{}\allowbreak{}jacobian-{}\allowbreak{}field-{}\allowbreak{}extension})
\end{description}
\clearpage\phantomsection%
\section*{References}
\addcontentsline{toc}{section}{References}
\bpbibitem{BBD82}{bp52}{A. Beilinson, J. Bernstein, and P. Deligne. \emph{Faisceaux pervers}, volume 100 of \emph{Astérisque}. 1982.}

\bpbibitem{BS15}{bp53}{B. Bhatt and P. Scholze. The pro-étale topology for schemes. \emph{Astérisque}, 369:99–201, 2015. \href{https://people.mpim-bonn.mpg.de/scholze/proetale.pdf}{Author's manuscript}.}

\bpbibitem{Con07}{bp54}{B. Conrad. Deligne's notes on Nagata compactifications. \emph{Journal of the Ramanujan Mathematical Society}, 22:205–257, 2007. \href{https://math.stanford.edu/\~conrad/papers/nagatafinal.pdf}{Author's manuscript}.}

\bpbibitem{Del74}{bp55}{P. Deligne. La conjecture de Weil. I. \emph{Publ. Math. IHÉS}, 43:273–307, 1974.}

\bpbibitem{Del80}{bp56}{P. Deligne. La conjecture de Weil. II. \emph{Publ. Math. IHÉS}, 52:137–252, 1980.}

\bpbibitem{Den22}{bp57}{C. Deninger. There is no “Weil-”cohomology theory with real coefficients for arithmetic curves. \emph{Ann. Sc. Norm. Super. Pisa Cl. Sci.}, 2022. To appear. arXiv:\href{https://arxiv.org/abs/2204.02714}{2204.02714}.}

\bpbibitem{Deu41}{bp58}{M. Deuring. Die Typen der Multiplikatorenringe elliptischer Funktionenkörper. \emph{Abh. Math. Sem. Univ. Hamburg}, 14:197–272, 1941.}

\bpbibitem{Eke90}{bp59}{T. Ekedahl. On the adic formalism. In \emph{The Grothendieck Festschrift, Vol. II}, volume 87 of \emph{Progress in Mathematics}, pages 197–218. Birkhäuser, 1990.}

\bpbibitem{Has36}{bp60}{H. Hasse. Zur Theorie der abstrakten elliptischen Funktionenkörper. I, II, III. \emph{J. Reine Angew. Math.}, 175:55–62, 69–88, 193–208, 1936.}

\bpbibitem{Hum72}{bp61}{J. E. Humphreys. \emph{Introduction to Lie Algebras and Representation Theory}, volume 9 of \emph{Graduate Texts in Mathematics}. Springer, 1972.}

\bpbibitem{Ill94}{bp62}{L. Illusie. Autour du théorème de monodromie locale. \emph{Astérisque}, 223:9–57, 1994. §4. \href{https://www.numdam.org/item/AST_1994__223__9_0/}{Original article}.}

\bpbibitem{Kat88}{bp63}{N. M. Katz. Travaux de Laumon. \emph{Astérisque}, 161–162:105–132, 1988. Séminaire Bourbaki, exp. 691.}

\bpbibitem{KW01}{bp64}{R. Kiehl and R. Weissauer. \emph{Weil conjectures, perverse sheaves and \(\ell\)-adic Fourier transform}, volume 42 of \emph{Ergebnisse der Mathematik}. Springer, 2001.}

\bpbibitem{Lau87}{bp65}{G. Laumon. Transformation de Fourier, constantes d'équations fonctionnelles et conjecture de Weil. \emph{Publ. Math. IHÉS}, 65:131–210, 1987.}

\bpbibitem{Mil86a}{bp66}{J. S. Milne. \href{https://www.jmilne.org/math/xnotes/AVs.pdf}{Abelian Varieties}. In \emph{Arithmetic Geometry}. 1986. Corrected online version.}

\bpbibitem{Mil86b}{bp67}{J. S. Milne. \href{https://www.jmilne.org/math/xnotes/JVs.pdf}{Jacobian Varieties}. In \emph{Arithmetic Geometry}. 1986. Corrected online version.}

\bpbibitem{Ran39}{bp68}{R. A. Rankin. Contributions to the theory of Ramanujan's function \(\tau(n)\) and similar arithmetical functions. II. \emph{Proc. Cambridge Philos. Soc.}, 35:357–372, 1939.}

\bpbibitem{SGA4}{bp69}{M. Artin, A. Grothendieck, and J.-L. Verdier, editors. \emph{Théorie des topos et cohomologie étale des schémas}, volumes 269, 270, 305 of \emph{Lecture Notes in Mathematics}. Springer, 1972–1973. SGA 4.}

\bpbibitem{SGA4½}{bp70}{P. Deligne. \emph{Cohomologie étale}, volume 569 of \emph{Lecture Notes in Mathematics}. Springer, 1977. SGA 4½; \emph{Théorèmes de finitude}, §3, and \emph{Rapport sur la formule des traces}, §3, and \emph{Sommes trigonométriques}, §1.}

\bpbibitem{SGA5}{bp71}{A. Grothendieck, editor. \emph{Cohomologie \(\ell\)-adique et fonctions \(L\)}, volume 589 of \emph{Lecture Notes in Mathematics}. Springer, 1977. SGA 5.}

\bpbibitem{SGA7II}{bp72}{P. Deligne and N. Katz, editors. \emph{Groupes de monodromie en géométrie algébrique}, volume 340 of \emph{Lecture Notes in Mathematics}. Springer, 1973. SGA 7 II; Exposé XIII, (2.1.7).}

\bpbibitem{Stacks}{bp73}{The Stacks Project Authors. \href{https://stacks.math.columbia.edu}{The Stacks Project}. Individual results are cited by their stable tags in the text.}

\bpbibitem{Var05}{bp74}{Y. Varshavsky. Lefschetz–Verdier trace formula and a generalization of a theorem of Fujiwara. 2005. Version 2. arXiv:\href{https://arxiv.org/abs/math/0505564}{math/0505564}. \href{https://arxiv.org/abs/math/0505564v2}{Author's preprint}.}

\bpbibitem{Wei48a}{bp75}{A. Weil. \emph{Sur les courbes algébriques et les variétés qui s'en déduisent}. Hermann, Paris, 1948.}

\bpbibitem{Wei48b}{bp76}{A. Weil. \emph{Variétés abéliennes et courbes algébriques}. Hermann, Paris, 1948.}

\bpbibitem{Wei49}{bp77}{A. Weil. Numbers of solutions of equations in finite fields. \emph{Bull. Amer. Math. Soc.}, 55:497–508, 1949.}

\end{document}
